% AVOI -- Stage 3 manuscript-ready results text (conservative). Numbers from results/stage3/tables.
% Stage 3 = adaptive next-evidence selection (WP1) and external-domain (AIT-LDS v2.0, CERT) evidence-model validation (WP2).
% This section reports Stage 3 on its own terms; it was designed and frozen after Stage 2 concluded and was not part
% of the original Stage-1 pre-registration or the Stage-2 protocol.
\section{Results: Stage 3 (adaptive evidence selection and external-domain validation)}

\subsection{Design}
Stage 3 was frozen (\texttt{config/stage3.yaml}) after Stage 2 concluded, using all nine Stage-1/2 evidence-model
datasets only as Stage-3 development data. Two questions were pre-registered. WP1: does AVOI's next-evidence choice
depend on previously observed evidence (A1), and does re-ranking after each observation reduce expected loss relative
to a policy that commits to a ranking computed once and never revised, at matched evidence budgets (A2, ``Frozen-Order
AVOI'')? WP2: does at least one external (non-IAM) evidence model pass the unchanged Stage-2 adequacy policy on
independent data (E1), and does AVOI retain an acquisition-efficiency advantage on it (E2)? The AVOI decision rule,
loss matrix, and adequacy policy are unchanged from Stage 2. New IAM scenarios use a scenario generator extended with
a latent, AVOI-invisible ``branching mechanism'' so that the same admissible evidence-action set can require different
next evidence depending on history. New external data are three previously unused AIT-LDS v2.0 testbeds and two
previously unused CERT releases.

\subsection{Adaptive next-evidence selection (A1)}
Across three IAM objectives with branching mechanisms, a group-aware permutation test (null: next action independent
of observed history, given objective and step; history permuted within scenario-family$\times$step blocks; 10{,}000
permutations) rejected independence for every objective (Holm-adjusted $p=0.0003$, all three). At least two objectives
selected more than one distinct next action at step $\ge 2$ across confirmation episodes, and rank-reversal rates
(the adaptive next action differing from the policy's own frozen initial ranking) were 55--59\% of observed steps.
AVOI's next-evidence choice is therefore genuinely history-dependent, not an artefact of a fixed initial ranking.

\subsection{Adaptive versus Frozen-Order AVOI (A2)}
Frozen-Order AVOI computes the AVOI ranking once, at the first decision point, and always acquires the
highest-ranked remaining action from that fixed list; every other component is identical to Adaptive AVOI. Both
policies were forced to acquire exactly $k\in\{2,3\}$ actions. Pooled over the three objectives (as pre-registered),
neither $k$ showed a 95\% paired-bootstrap confidence interval on the loss difference entirely below zero ($k=2$:
$-0.022$, CI $[-0.048,+0.010]$; $k=3$: $+0.003$, CI $[-0.038,+0.048]$), and neither reached Holm significance
(sign-flip $p=0.143$ and $0.951$). \textbf{A2 is not supported at the pre-registered pooled level.} Reported
descriptively per objective: on IAM-02-TERM-DISABLED, adaptive re-ranking was significantly better than the frozen
order at both horizons (loss lower by 0.086--0.113, Holm $p=0.0006$). On IAM-01-MFA-PRIV the two policies were
outcome-identical (the acquired evidence set, not its order, determines the terminal decision on this objective). On
IAM-03-PRIV-AUTHZ -- the one WP1 objective whose evidence model failed the adequacy calibration condition -- the
frozen order was numerically better, reaching a marginal effect at $k=3$ (Holm $p=0.098$): adaptive re-ranking against
a poorly calibrated value signal can cost loss rather than save it.

\subsection{External evidence-model adequacy (E1)}
The Stage-2 AIT exfiltration-detection failure (AUROC 0.024) was diagnosed as a non-monotone calibration-map artefact
of a flawed two-testbed development split, not a genuine inversion: the raw discriminating statistic separates
violations with AUROC 1.00 on every development testbed examined. A calibrator restricted to monotone maps, combined
with leave-one-development-testbed-out model selection, removes this failure mode by construction. On three
previously unused AIT-LDS v2.0 testbeds and two previously unused CERT releases, one of three external claims passed
the unchanged adequacy policy: CERT-INSIDER (AUROC 0.895, 95\% CI $[0.671,0.953]$; Brier skill score 0.744, CI
$[0.554,0.918]$; expected calibration error 0.017; $n=394$, 21 violated episodes, 187 independent groups) was
VALIDATED under all primary adequacy conditions. Both AIT claims remained RESTRICTED for insufficient sample support
on the pre-committed target testbed; one of the two (intrusion detection) showed strong, robust discrimination
(worst-development-fold AUROC 0.837; target-testbed AUROC 0.994; confirmation-testbed AUROC 0.969) and would plausibly
validate under a target assignment with more positive cases, while the other (exfiltration detection) showed weaker,
less consistent transfer (worst-development-fold AUROC 0.638; confirmation-testbed AUROC 0.729, falling to 0.563 on
one confirmation testbed) independent of its sample-support failure.

CERT-INSIDER's pass is not evidence of good operating-point performance: at its adequacy-validation-tuned decision
threshold, sensitivity to the violated class was 0.0 (0 of 21 violated episodes correctly flagged) despite good ranking
and calibration quality. The adequacy policy's conditions concern probabilistic discrimination, calibration and sample
support, not threshold-level recall of the class of interest.

\subsection{Acquisition efficiency on the validated external model (E2)}
On CERT-INSIDER confirmation episodes ($n=458$), AVOI acquired 2.42 items on average versus 4.85 for a fixed checklist
and 10.0 for exhaustive collection (paired Wilcoxon, Holm-adjusted $p<10^{-75}$ throughout), at normalised cost 0.056
versus 0.416 and 1.237 respectively. Applying a pre-registered no-trivial-abstention rule (efficiency is claimed only
where AVOI's autonomous coverage is not below the comparator's by more than 0.10): the efficiency advantage holds
against the fixed checklist (coverage gap 0.026) but is not claimed against exhaustive collection (coverage gap 0.150)
or against information gain, whose operating point on this domain was degenerate (zero acquisitions, 100\% false
assurance among its ``assured'' decisions). \textbf{E2 is therefore partially supported}: an efficiency advantage
exists, but only relative to a subset of comparators, on the single external claim that reached VALIDATED status.

\subsection{Limitations}
The CERT result uses a release pool reduced from four to two releases (7 pooled insiders rather than the originally
planned 17), a time-budget decision made and frozen before either new release's label was read; the two remaining
releases were fully downloaded and verified but not processed. Both external-data-independence structures (the
size-based AIT target-testbed assignment; the reduced CERT release pool) are single, non-adversarial choices, not
robustness sweeps. A2's pooled test combines one objective with a significant adaptive win, one with a mechanically
null result, and one with an adequacy-limited loss, into a single pre-registered statistic; this is reported as
designed, not re-aggregated post hoc to change the conclusion.
